\documentclass[12pt,openany]{extbook}
% \documentclass[a4paper,17pt]{extarticle}
\usepackage[a4paper, total={7in, 10in}]{geometry}
\usepackage[utf8]{inputenc}
\DeclareUnicodeCharacter{2212}{-}
\DeclareUnicodeCharacter{2502}{|}
\DeclareUnicodeCharacter{200B}{ }
\usepackage[russian]{babel}
\usepackage{graphicx}
\usepackage{float}
\usepackage{stackengine}
\usepackage{multicol}
\usepackage{amsfonts}
\usepackage{hyperref}
\usepackage[margin=10pt,font=footnotesize,labelfont=bf,labelsep=endash]{caption}
\newenvironment{Figure}{\par\medskip\noindent\minipage{\linewidth}}{\endminipage\par\medskip}
\setlength{\headheight}{14.5pt}
\usepackage{titling}
\newcommand{\subtitle}[1]{%
  \posttitle{%
    \par\end{center}
    \begin{center}\large#1\end{center}
   }%
}
\renewcommand{\baselinestretch}{1.15} 


% Для химических формул
% \usepackage[version=4]{mhchem}

% Для матеши
% \usepackage{amsmath}
% \usepackage{amssymb}

\usepackage{fancyhdr}
\newcommand{\bosversion}{v1.0.\hbox{\today}}

\pagestyle{fancy}
\fancyhf{}
\rhead{\rightmark}
\lhead{\leftmark}
\makeatletter
\let\ps@plain\ps@fancy % plain style = fancy style
\makeatother
% \rfoot{Стр. \thepage}
\fancyfoot[LE,RO]{Стр. \thepage/24}
\fancyfoot[RE,LO]{qazcho.kz}
% \fancyfoot[R]{\thepage} % except the right
% \renewcommand{\contentsname}{Hello}
% \addto\captionrussian{% Replace "english" with the language you use
%   \renewcommand{\contentsname}%
%     {Hello}%
% }
% \usepackage{helvet}

% \renewcommand{\sfdefault}{ptm}
% \usepackage{sectsty}
% \allsectionsfont{\sffamily\bfseries}
% \setlength{\beforechapskip}{20pt}
\makeatletter
\renewcommand{\@chapapp}{Раздел}
\makeatother

\usepackage[x11names]{xcolor}
\usepackage{framed}
\usepackage{quoting}

 \colorlet{shadecolor}{LavenderBlush2}
\usepackage{lipsum}
\newenvironment{shadedquotation}
 {\begin{shaded*}
  \quoting[leftmargin=0pt, vskip=0pt]
 }
 {\endquoting
 \end{shaded*}
}

\usepackage{xcolor}

\usepackage{lipsum} % for mock text

\newsavebox{\coloredquotationbox}
\newenvironment{coloredquotation}
 {%
  \begin{trivlist}
  \begin{lrbox}{\coloredquotationbox}
  \begin{minipage}{\dimexpr\linewidth-2\fboxsep}
 }
 {%
  \end{minipage}
  \end{lrbox}
  \item\relax
  \parbox{\linewidth}{
    \begingroup
    \color[RGB]{224,215,188}%
    \hrule
    \color[RGB]{249,245,233}%
    \hrule
    \color[RGB]{224,215,188}%
    \hrule
    \endgroup
    \colorbox[RGB]{249,245,233}{\usebox{\coloredquotationbox}}\par\nointerlineskip
    \begingroup
    \color[RGB]{224,215,188}%
    \hrule
    \color[RGB]{249,245,233}%
    \hrule
    \color[RGB]{224,215,188}%
    \hrule
    \endgroup
  }
  \end{trivlist}
 }

\usepackage{titlesec}

\titlespacing*{\section} {0pt}{2.5ex plus 1ex minus .2ex}{1.3ex plus .2ex}
\titlespacing*{\subsection} {0pt}{2.25ex plus 1ex minus .2ex}{0.5ex plus .2ex}

\begin{document}

\begin{titlepage}


\newcommand{\HRule}{\rule{\linewidth}{0.5mm}} % Defines a new command for the horizontal lines, change thickness here

\center % Center everything on the page

\includegraphics[width=0.5\linewidth]{logo_color.jpg}\\[1.0cm]

%----------------------------------------------------------------------------------------
%	HEADING SECTIONS
%----------------------------------------------------------------------------------------

\begin{minipage}{0.5\textwidth}
\begin{flushleft} \large
\textsc{\textbf{Республиканская Коллегия Составителей Задач по Олимпиадной Химии}}
\end{flushleft}
\end{minipage}
~
\begin{minipage}{0.4\textwidth}
\begin{flushright} \large
\textsc{\textbf{Профиль участника РО (2022-2023 уч. год)}}\\
\textsc{\today}
\end{flushright}
\end{minipage}\\[2.0cm]

\begin{minipage}{0.93\textwidth}
\begin{flushleft}
  

\textbf{\large Введение}\\[0.5cm]
В 2022-2023 уч. году, комплекты заданий районного, областного и заключительного этапов республиканской олимпиады школьников по химии (далее - РО) сопровождались призывом к участникам оставить обратную связь по прошедшей олимпиаде путем заполнения онлайн опроса, созданного Республиканской Коллегией Составителей Задач по Олимпиадной Химии (далее - Коллегия). 

Первостепенной задачей опроса было получение информации о том, как ученики готовятся к олимпиадам и насколько сложными, объемными и интересными им показались задания. Дополнительно была собрана демографическая информация. В этой заметке предоставляем результат анализа собранных ответов. Главная задача при написании этого документа - улучшить представления составителей заданий о том, какой уровень подготовки имеется у участников того или иного этапа РО. 
\end{flushleft}
\end{minipage}


\end{titlepage}


\def\contentsname{Содержание}
\tableofcontents

\chapter{Демографический профиль участника РО}
\section{Введение}

В 2022-2023 учебном году, комплекты заданий и решений теоретических туров районного, областного и заключительного этапов республиканской олимпиады школьников по химии сопровождались призывом оставить обратную связь о прошедшей олимпиаде, заполнив онлайн опрос, созданный Коллегией. Ссылка на опрос также дублировалась на официальных каналах Коллегии и на (\href{https://ask.bc-pf.org}{форуме Спроси}). В распространении опроса РНПЦ Дарын и Министерство Просвещения РК не принимали никакого участия. В рамках опроса, помимо «существенных» вопросов, ответы на которые проанализированы в других разделах, участники предоставляли демографические данные. Некоторые сводки настолько нас удивили, что мы решили посвятить им отдельную секцию.

Опрос по районному этапу РО прошли 40 участников. Из них 10 (25\%), 15 (37.5\%) и 15 (37.5\%) человек участвовали за 9, 10 и 11 класс соответственно. Распределение по школам следующее:

\begin{table}[ht]
    \centering
    \begin{tabular}{|l|c|c|}
    \hline
    \textbf{Тип школы} & \textbf{Количество} & \textbf{Процент} \\ \hline
    Обычная школа (в городе) & 16 & 40.0\% \\ \hline
    Обычная школа (в селе) & 12 & 30.0\% \\ \hline
    БИЛ & 8 & 20.0\% \\ \hline
    Другая специализированная школа & 2 & 5.0\% \\ \hline
    Частная школа & 2 & 5.0\% \\ \hline
    НИШ & 0 & 0.0\% \\ \hline
    РФМШ & 0 & 0.0\% \\ \hline
    \end{tabular}
    \caption{Типы школ и их распределение на районном этапе РО}
    \label{tab:region-schols}
\end{table}

Стоит заметить, что ученики НИШ и БИЛ не пишут районный этап и проходят сразу на областной, поэтому ответов от учеников НИШ нет. 8 ответов от учеников БИЛ скорее всего предоставлены учениками, которые самостоятельно прорешивали задания районного этапа.

Опрос по областному этапу РО прошли 60 участников. Из них 28 (46.7\%), 17 (28.3\%), и 15 (25.0\%) человек участвовали за 9, 10 и 11 класс соответственно. Распределение по школам следующее:

\begin{table}[ht]
    \centering
    \begin{tabular}{|l|c|c|}
        \hline
        \textbf{Тип школы} & \textbf{Количество} & \textbf{Процент} \\ \hline
        БИЛ & 31 & 51.7\% \\ \hline
        Другая специализированная школа & 11 & 18.3\% \\ \hline
        НИШ & 8 & 13.3\% \\ \hline
        Обычная школа (в городе) & 6 & 10.0\% \\ \hline
        Частная школа & 3 & 5.0\% \\ \hline
        РФМШ & 1 & 1.7\% \\ \hline
        Обычная школа (в селе) & 0 & 0.0\% \\ \hline
    \end{tabular}
    \caption{Типы школ и их распределение на областном этапе РО}
    \label{tab:obl-schools}
\end{table}

Нельзя не заметить практически полное отсутствие ответов от учников из неспециализированных школ. Конечно, 40 и 60 респондентов для этапа, на котором участвуют десятки тысяч (район) и тысячи (область) школьников, могут быть не репрезентативной выборкой. Однако, сложно представить, что ученики из неспециализированных школ, которые участвовали на районном этапе и заполняли опрос, не заполнили бы опрос на областном этапе, если бы они на нем участвовали. Таким образом, Таблица \ref{tab:obl-schools} может указывать на то, что ученики неспециализированных школ в своем большинстве не проходят дальше районного этапа.

Опрос по заключительному этапу РО прошли 32 участника. Из них 9 (28.1\%), 9 (28.1\%) и 14 (43.8\%) человек участвовали за 9, 10 и 11 класс соответственно. Распределение по школам следующее:

\begin{table}[ht]
    \centering
    \begin{tabular}{|l|c|c|}
    \hline
    \textbf{Тип школы} & \textbf{Количество} & \textbf{Процент} \\ \hline
    БИЛ & 15 & 46.9\% \\ \hline
    Другая специализированная школа & 5 & 15.6\% \\ \hline
    Обычная школа (в городе) & 5 & 15.6\% \\ \hline
    НИШ & 4 & 12.5\% \\ \hline
    Частная школа & 2 & 6.3\% \\ \hline
    РФМШ & 1 & 3.1\% \\ \hline
    Обычная школа (в селе) & 0 & 0.0\% \\ \hline
    \end{tabular}
    \caption{Типы школ и их распределение на заключительном этапе РО}
    \label{tab:respa-schools}
\end{table}

\section{Язык обучения и язык участия на РО}

% \input{tex/participants}

% \chapter{Вопросы логистики}
% \input{tex/logistics}

% \chapter{Вопросы подготовки заданий}
% \input{tex/problemprep}

% \chapter{Приложения}
% \input{tex/appendix}


\end{document}
